% Edited conversion of source pp. 28--29.
\ReportHeading
  {Застосування асимптотичного методу Павленка--Маневича у сучасних контактних задачах механіки}
  {Т.~С. Кагадій\textsuperscript{1}, О.~В. Білова\textsuperscript{2}, І.~В. Щербина\textsuperscript{3}, А.~Г. Шпорта\textsuperscript{1}}
  {\textsuperscript{1}Національний технічний університет ``Дніпровська політехніка''; \textsuperscript{2}Український державний університет науки і технологій; \textsuperscript{3}Дніпровський державний аграрно-економічний університет}
  {}

Асимптотичний метод Павленка--Маневича є ефективним підходом до розв’язування крайових задач механіки деформівного твердого тіла. Він ґрунтується на побудові асимптотичних розкладів шуканих величин за малим параметром, який характеризує геометричні або фізичні особливості системи. Такий підхід дає змогу виділити головні члени розкладу, що визначають напружено-деформований стан, і отримати аналітичні залежності для складних задач, які важко або неможливо дослідити прямими методами.

Суть методу полягає у зведенні вихідної задачі до послідовності задач нижчої розмірності з подальшим узгодженням їхніх розв’язків. Це забезпечує ефективне врахування локальних особливостей --- тонких включень, стрингерів, контактних зон або областей із різко змінними властивостями матеріалу --- без повного чисельного моделювання всієї області~\cite{pavlenko:kagadiy2016}.

Актуальність методу зростає у зв’язку з розвитком сучасних матеріалів, зокрема композитів із криволінійною анізотропією та електропружних, або п’єзоелектричних, середовищ. У таких системах класичні підходи часто є неефективними через складну взаємодію механічних і фізичних полів. Асимптотичний метод дає змогу врахувати ці ефекти шляхом введення відповідних малих параметрів і побудови узгоджених розкладів для механічних та електричних величин~\cite{pavlenko:kagadiy2016,pavlenko:kagadiy2021}.

Особливе значення метод має для контактних задач, у яких локальні напруження в областях контакту визначають міцність і довговічність конструкцій. Асимптотичні розклади дають аналітичні вирази для контактних напружень і дають змогу оцінювати вплив анізотропії, геометрії та фізичних властивостей матеріалу на розподіл полів і встановлювати основні закономірності поведінки системи. Метод також ефективний для конструкцій із тонкими елементами та кількома масштабами довжин: він зменшує обчислювальні витрати, зберігаючи фізичну наочність результатів.

Отже, асимптотичний метод Павленка--Маневича поєднує аналітичну точність із можливістю дослідження складних багатопараметричних систем, зокрема контактних задач для композитних і електропружних матеріалів.

\begin{thebibliography}{9}
\bibitem{pavlenko:kagadiy2016}
Т.~С. Кагадій, О.~В. Білова, І.~В. Щербина,
``Аналітичний підхід до розв’язання деяких контактних задач,''
\emph{Вісник Херсонського національного університету}, №~3(58), с.~104--110, 2016.

\bibitem{pavlenko:kagadiy2021}
Т.~С. Кагадій, А.~Г. Шпорта, О.~В. Білова, І.~В. Щербина,
``Математичне моделювання в задачах геометрично нелінійної теорії пружності,''
\emph{Прикладні питання математичного моделювання}, т.~4, №~1, с.~103--110, 2021.
\url{https://doi.org/10.32782/KNTU2618-0340/2021.4.1.11}.
\end{thebibliography}
